%% kapa: exam=1 round=35 year=2024 period=2 form=A subject=appraisal_law subject_ko=감정평가관계법규 seq=21 number=21
%% source: Q-Net 감정평가사 기출문제 (cst00309) artlSeq=5213229 file=raw/1008/5213229/2교시 문제지 원본.pdf
%% conversion: kordoc 4.24.51 -> tools/gen1.py (auto); review=proofread-sonnet; flags=none
\begin{question}{21}
국유재산법령상 국유재산에 관한 설명으로 옳지 않은 것은?
\begin{choices}
  \item 정부시설에서 사용하는 궤도차량으로서 해당 시설의 폐지와 함께 포괄적으로 용도폐지된 것은 해당 시설이 폐지된 후에는 국유재산으로 하지 아니한다.
  \item 총괄청은 일반재산을 보존용재산으로 전환하여 관리할 수 있다.
  \item 등기가 필요한 국유재산이 부동산인 경우 그 권리자의 명의는 국(國)으로 하되 소관 중앙관서의 명칭을 함께 적어야 한다.
  \item 총괄청이나 중앙관서의 장은 소유자 없는 부동산을 국유재산으로 취득한다.
  \item 지상권, 전세권, 광업권은 국유재산의 범위에 속한다.
\end{choices}
\end{question}
\answer{1}
